\section{A Gaussian check of the blocked-sampler comparison}\label{sec:gausscheck}

Section~\mref{sec:blocked} states two things about $\alpha$: that the
lag-1 autocorrelation of an exact two-block sampler equals $R^2$ for a
Gaussian posterior (Proposition~\mref{prop:blocked}(iii)), and,
heuristically, that HMC pays for the same correlation at most in
integration steps. We checked both in a Gaussian model, comparing a
simulated two-block Gibbs sampler with a stylized model of HMC (exact
Hamiltonian dynamics, random integration time). The two-block sampler's
lag-1 autocorrelation matched $R^2$ to within 0.004, and its integrated
autocorrelation time went 3.0, 9.0, 27.6, 39.0 and 199 as $R^2$ went 0.5,
0.8, 0.93, 0.95 and 0.99. The HMC model's stayed between 1.6 and 2.1, at a
cost of 4.8 to 46.8 gradient evaluations per independent draw. With an
unrelated narrow ridge already present (such as, in a joint model, the
spline coefficients or variance components might provide), the step count was 22.2 at
every $R^2$ from 0 to 0.99.

\section{Random intercept only: runs to convergence}\label{sec:longrun}

The same simulator as in Section~\mref{sec:simulation} with $\sigma_1 =
0$, fitted with a random intercept only, gave seven designs: about 9, 17
and 35 visits, $N \in \{100, 300\}$, and one noisier design with $\sigma_e
= 0.8$. The intercept gained 6.5--282 times and the subject-constant
coefficient of age 6.6--207 times; the time slope, which has no absorbable
direction when there is no random slope, moved only through the step size
(0.79--3.4 times). At this run length, however, 17 of the 21 unrotated fits
had some parameter with $\hat R$ above 1.05, the intercept's reaching 1.27,
while every rotated fit had $\hat R \leq 1.004$, so these ratios compare a
converged fit with unconverged ones. Posterior means agreed to a median of
0.025 posterior standard deviations (maximum 0.175).

To obtain a comparison between converged fits, we reran three designs
($N = 300$ with about 9, 17 and 35 visits) and extended the unrotated fit
to 2, 4, 8, 16 and 32 times the standard length, with two seeds taken to
the longest runs. A convergence criterion was fixed in advance: split $\hat
R \leq 1.01$ on every monitored parameter and ESS $\geq 400$ for the
intercept, age, the time slope and $\alpha$. Every rotated fit met it at
the standard length. The unrotated fit met it first at 16 times the
standard length with 9 visits, at 16 and 32 times with 17 visits, and not
by 32 times with 35 visits, where its intercept had an ESS of 304--377 from
64,000 draws with $\hat R = 1.002$. Its ESS per draw for the intercept did
not change with run length: 0.009--0.023, 0.007--0.020 and 0.002--0.008 in
the three designs over runs from 1 to 32 times the standard length, against
2.1--2.5, 0.9--1.9 and 1.1--2.4 for the rotated fit. The per-draw advantage,
of order 100 times, is therefore a property of the stationary chains, not
of short runs. In wall time, the unrotated fit needed 3.5 and 6.5--6.6
times as long as the rotated fit to meet the criterion in the first two
designs, and more than 10--14 times in the third. These ratios are much
smaller than the per-draw ones because the rotated fit's time at the
standard length is mostly warm-up and compilation, and it exceeded the ESS
target by a factor of 5 to 12. Between the first fits of each arm to meet
the criterion, posterior means agreed to a median of 0.02 posterior
standard deviations (maximum 0.11) and standard deviations to within 6\%.

\section{Stress grid}\label{sec:stressgrid}

The stress grid of Section~\mref{sec:stress} crossed about 35 visits with
$N \in \{50, 100\}$ and $\rho \in \{0.3, 0.95\}$, and added $\rho = 0.95$
at 9 visits ($N = 100$) and at 17 visits ($N = 300$). The criterion fixed before the
run was that the rotation must not fall below the unrotated fit on the
intercept or time slope in any design. It did not (Table~\ref{tab:stress});
the lowest design was 5.5 times and the lowest single seed 3.2 times. Per
gradient evaluation the gains were 16.7 for the intercept and 22.0 for the
time slope (geometric means over all designs), and the rotated fit took
0.81 times as many leapfrog steps per iteration.

\begin{table}[ht]
\centering
\caption{Stress grid: ESS-per-second ratios, rotated over unrotated
(geometric means over 3 seeds).}
\label{tab:stress}
\small
\begin{tabular}{rrr rrr}
\toprule
visits & $N$ & $\rho$ & intercept & time slope & $\alpha$ \\
\midrule
35 &  50 & 0.30 &  5.5 & 10.3 & 1.60 \\
35 &  50 & 0.95 & 21.1 & 21.6 & 2.42 \\
35 & 100 & 0.30 & 14.5 & 27.1 & 1.58 \\
35 & 100 & 0.95 & 26.0 & 28.6 & 1.73 \\
 9 & 100 & 0.95 &  7.0 &  7.5 & 0.74 \\
17 & 300 & 0.95 & 24.6 & 29.7 & 1.09 \\
\bottomrule
\end{tabular}
\end{table}

The design with the fewest subjects and the most information per subject
gained least, as the fixed-metric residual predicts. Strong correlation
did not reduce the rotation's advantage: at $\rho = 0.95$ the gains were
as large as or larger than at $\rho = 0.3$, partly because the unrotated
fit degrades faster. Nor did it hurt the rotated fit in absolute terms:
with 35 visits its ESS per draw for the intercept was 1.25 at $\rho = 0.95$
against 0.68 at $\rho = 0.3$ with $N = 50$, and 1.50 against 1.18 with $N
= 100$. The predicted weakness is visible only against a variant that
removes the absorbable directions from the likelihood's coordinates and
restores the reported coefficients with an exact correction
(Section~\mref{sec:related}). That variant has no fixed-metric residual,
and in the 35-visit designs it was 1.4--2 times more efficient than the
rotation, while being about equal to it with 9 visits. Posterior means
agreed with the unrotated fit's to a median of 0.029 posterior standard
deviations (maximum 0.166, 90 pairs), there were no divergences, and wall
time was 0.86--1.08 times the unrotated fit's (geometric means per
design). $\alpha$, which the degeneracy does not involve, fell to 0.74 in
the 9-visit design.

\section{Calibration}\label{sec:calib}

With $N = 300$ and 100 replicates, every one of which passed an $\hat R
\leq 1.05$ check on the regression parameters (so none was excluded),
coverage of 95\% intervals agreed between the unrotated and rotated fits
within 0.01, and the ratio of mean posterior standard deviation to the
standard deviation of posterior means within 0.015, for every estimand
(Table~\ref{tab:cal}); the Monte Carlo standard error of a 95\% coverage is
about 0.011. Two features appear in both arms and belong to the model and
design: slightly wide intervals for $\beta_0$, and an upward bias of about
3 Monte Carlo standard errors (2.8\%) in $\alpha$. The latter is consistent
with ordinary finite-sample bias. In a separate study
(\texttt{dev/study\_alpha\_bias.R}, Section~\ref{sec:repro}), neither the spline baseline hazard nor
the simulator's hazard calibration explained it (paired changes of
$-0.0011$ and $-0.0016$, standard error about 0.0025), and at $N = 1200$ it
fell to 1.1\%, 1.7 Monte Carlo standard errors, with coverage 0.92.

\begin{table}[ht]
\centering
\caption{Calibration, $N = 300$, 100 replicates: coverage of 95\%
intervals, ratio of mean posterior SD to the SD of posterior means, and
bias in Monte Carlo standard errors.}
\label{tab:cal}
\small
\begin{tabular}{l rrr rrr}
\toprule
& \multicolumn{3}{c}{unrotated} & \multicolumn{3}{c}{rotated} \\
\cmidrule(lr){2-4}\cmidrule(lr){5-7}
estimand & cover & SD ratio & bias/se & cover & SD ratio & bias/se \\
\midrule
$\beta_0$           & 0.99 & 1.164 &  1.68 & 0.99 & 1.150 &  1.55 \\
$\beta_{\mathrm{time}}$ & 0.97 & 1.027 & $-0.30$ & 0.96 & 1.038 & $-0.32$ \\
$\alpha$            & 0.93 & 0.922 &  2.97 & 0.93 & 0.917 &  2.97 \\
$\sigma_0$          & 0.94 & 0.914 &  0.63 & 0.95 & 0.906 &  0.68 \\
\bottomrule
\end{tabular}
\end{table}

\section{Real data with a random intercept only}\label{sec:realq1}

Fitted with a random intercept only, the two datasets of
Section~\mref{sec:realdata} show the degeneracy more starkly. Unrotated,
$\beta_0$ and the subject-constant coefficients mixed at 0.04--0.07 ESS per
draw, and 91--95\% of their posterior variance was explained by the
absorbable directions of the random effects; every other parameter mixed
near one ESS per draw. The rotation raised the affected coefficients' ESS
per second 51--53 times on \texttt{aids} and 19--20 times on \texttt{pbc2}
(2 seeds each), and the slowest parameter's ESS per draw from 0.05 to 0.88
and from 0.06 to 0.78.

\section{Reproducibility}\label{sec:repro}

The scripts below are in the \texttt{dev/} folder of the \texttt{jmjax}
source at release 0.3.0 (DOI 10.5281/zenodo.22950084;
\url{https://github.com/jmjaxmaintainer-jpg/jmjax}), and paths are
relative to its root. The article's code archive contains that source,
the two result files written outside the repository
(\texttt{results/}), and a README with setup and run instructions.

\begin{center}
\footnotesize
\begin{tabular}{>{\raggedright\arraybackslash}p{0.33\textwidth}>{\raggedright\arraybackslash}p{0.60\textwidth}}
\toprule
result & script (output) \\
\midrule
Figure~\mref{fig:ridge}, Section~\mref{sec:residual} numbers &
\texttt{dev/theory\_rotation\_toy.py}; \texttt{dev/figures/make\_figures.py} \\
Section~\ref{sec:gausscheck} & \texttt{dev/theory\_alpha\_toy.py} \\
Section~\mref{sec:blocked}, $R^2$ and its split & \texttt{dev/alpha\_missing\_information.R} \\
Figure~\mref{fig:alpha} & \texttt{dev/figures/make\_figures.py} \\
Table~\mref{tab:core} & \texttt{dev/pilot\_rotate\_grid.R} (\texttt{dev/pilot\_rotate\_grid.csv}) \\
Section~\mref{sec:q1} and Section~\ref{sec:longrun}, standard length & \texttt{ROT\_GRID=q1 dev/pilot\_rotate\_grid.R} (\texttt{dev/pilot\_rotate\_q1.csv}) \\
Section~\ref{sec:longrun} & \texttt{dev/study\_q1\_longrun.R} (\texttt{dev/study\_q1\_longrun.csv}) \\
Table~\ref{tab:stress} & \texttt{ROT\_GRID=stress dev/pilot\_rotate\_grid.R} (\texttt{dev/pilot\_rotate\_stress.csv}) \\
Figure~\mref{fig:gain} & \texttt{dev/figures/make\_figures.py} \\
Table~\ref{tab:cal} & \texttt{dev/study\_calibration.R} (\texttt{dev/calibration\_results.csv}) \\
$\alpha$ bias, $N = 1200$ & \texttt{dev/study\_alpha\_bias.R} (\texttt{dev/alpha\_bias\_results.csv}) \\
Table~\mref{tab:gen} & \texttt{dev/generality/mixed\_models.py} \\
& (\texttt{dev/generality/results\_orthodont\_toenail\_ext.csv}) \\
Table~\mref{tab:real} & \texttt{dev/study\_realdata\_rotate.R} (\texttt{results/study\_realdata\_rotate.csv}) \\
Table~\mref{tab:jmb}; JMbayes2 in Figure~\mref{fig:alpha} & \texttt{dev/study\_common\_ess.R} (\texttt{dev/common\_ess.csv}) \\
Section~\ref{sec:realq1} & \texttt{dev/pilot\_q1\_intercept.R} (\texttt{results/pilot\_q1\_intercept.csv}) \\
$\alpha$ against JMbayes2, 1--12 covariates & \texttt{dev/study\_orth\_vs\_jmbayes2.R} \\
Example~\mref{ex:spline}, basis checks & \texttt{dev/verify\_absorbable\_basis.py} \\
\bottomrule
\end{tabular}
\end{center}
